% Values: ILLUMINATING
% Evidence planned:
% - ILLUMINATING: how an implementation is organised (modules elaborated,
%   then evaluated, then the program run by the platform); the three phases
%   and the static and dynamic basis; the levels Core, Modules and Programs
%   and the Effects chapter; the grid table, cell for cell with the plan's
%   Grid; bare language and derived forms; one sentence per appendix; the
%   form of sentences; pointers to where the notational conventions are
%   defined (finite maps, projection, options, sequences) and the convention
%   for decorated metavariables.

\chapter{Introduction}\label{ch:introduction}

This book formally defines Standard Elm, the language Elm 0.19.1 as
delimited in the Preface.

The method of definition is easiest to understand by considering how an
implementation of Elm is naturally organised. A \emph{program} consists of
modules: a main module, and the modules it imports, directly or indirectly,
with no cycle among the imports. The implementation treats the modules one at
a time, each after the modules it imports. Each module is checked for the
consistency of its types and for other conditions of form, and what it makes available to other modules is
recorded in the top-level environment, which we call a \emph{basis}; the
modules that import it are then checked against that basis. When every
module has been checked, the top-level values of the modules are computed,
again each module after those it imports, and recorded in the basis. Finally
the platform runs the program: it starts it, performs the commands it
issues, maintains its subscriptions, and delivers to it, one at a time, the
messages that result.

Two words name the parties to this last step. The \emph{platform} is the
part of an implementation that runs a program, performs its effects and
exchanges values with the host. The \emph{host} is the environment outside
the program: a JavaScript page, an embedding application, or an operating
system. The host sends values to the incoming ports of a program and
receives those of its outgoing ports, through the platform. The later
chapters use both words with these meanings only. % decision 7 (audit Y11)

There are thus three phases: \emph{parsing}, \emph{elaboration} and
\emph{evaluation}. Parsing determines the grammatical form of a module.
Elaboration, the \emph{static} phase, determines whether its types are
consistent and whether it meets the other conditions of form; a module that
does both \emph{elaborates}. Elaboration records type and interface information in the
basis. Evaluation, the \emph{dynamic} phase, determines the values of its
declarations and records them in the basis. Corresponding to the phases,
the definition divides into grammatical rules, elaboration rules and
evaluation rules, and the basis divides into the \emph{static basis} and the
\emph{dynamic basis}: a top-level variable, for example, is associated with a
type scheme in the static basis and with a value in the dynamic basis. An
implementation need not divide the basis so; for the definition the
division makes each phase simpler to present. Every module of a program is
elaborated before any is evaluated, and a program that does not elaborate
is not evaluated at all.

A further factoring follows the structure of the language. Standard Elm has
a lower level, the \emph{Core}, of expressions, patterns, types and
\kw{let} definitions: everything the right-hand side of a single definition
can contain. Above it is \emph{Modules}: module headers, exposing lists,
imports, and the top-level declarations of a module, namely value
definitions, custom types, type aliases and ports. The top level,
\emph{Programs}, defines what a program is, the type its \kw{main} must
have, and how the platform runs it; all three phases of Programs are treated
in one chapter. Commands, subscriptions and tasks add no syntax, being types
of the basis, but their values and the behaviour of ports need a semantics
of their own. We give it in a chapter on \emph{Effects}, which is purely
dynamic; the static question of which types may cross a port belongs to the
static semantics of Modules. % D10, D11
The scheme of the chapters is therefore as follows:
\begin{center}
\setlength{\tabcolsep}{4pt}%
\begin{tabular}{r|c|c|c|c|}
  \multicolumn{1}{r}{}
  & \multicolumn{1}{c}{\emph{Core}}
  & \multicolumn{1}{c}{\emph{Modules}}
  & \multicolumn{1}{c}{\emph{Effects}}
  & \multicolumn{1}{c}{\emph{Programs}}\\
  \cline{2-5}
  \emph{Syntax}
  & Chapter~\ref{ch:syntax-core}
  & Chapter~\ref{ch:syntax-modules}
  & none
  & \\
  \cline{2-4}
  \emph{Static Semantics}
  & Chapter~\ref{ch:static-core}
  & Chapter~\ref{ch:static-modules}
  & Section~\ref{sec:static-modules-ports}
  & Chapter~\ref{ch:programs}\\
  \cline{2-4}
  \emph{Dynamic Semantics}
  & Chapter~\ref{ch:dynamic-core}
  & Chapter~\ref{ch:dynamic-modules}
  & Chapter~\ref{ch:effects}
  & \\
  \cline{2-5}
\end{tabular}
\end{center}

The Core and Modules provide many forms of phrase for the programmer's
convenience. Many of them are \emph{derived forms}, whose meaning is given by
translation into the remaining phrases, which we call the \emph{bare}
language: infixed operators, conditionals, \kw{\&\&} and \kw{||}, negation,
list expressions and patterns, field selection and field accessors,
functions of several arguments, and the grouping of the definitions of a
\kw{let} or of a module body into dependency groups. Each semantics chapter
treats only the bare language; the derived forms are given in
Appendix~\ref{ch:appendix-derived-forms}. The full grammar of the language,
bare and derived together, is collected in
Appendix~\ref{ch:appendix-full-grammar}.

Appendices~\ref{ch:appendix-static-basis} and~\ref{ch:appendix-dynamic-basis}
give the \emph{initial basis}, divided into its static and dynamic parts:
the predefined types, constructors and values that the rules and the
default imports need, with the basic arithmetic and comparison operations
defined there precisely. % D4
These operations are total, so that no evaluation ends in a run-time
error. % D2, D3
Where implementations legitimately differ, the body of the definition
states only what each must satisfy, and
Appendix~\ref{ch:appendix-implementation} records the choices of two
implementations and every known point on which they do not conform. % D8
Appendix~\ref{ch:appendix-standard-ml} compares this definition with the
Definition of Standard ML.

The semantics is presented as a set of rules allowing \emph{sentences} of the
form
\[
  \jdg{A}{\mathit{phrase}}{A'}
\]
to be inferred, where $A$ is usually a basis or a context (static) or an
environment (dynamic), and $A'$ is a semantic object, usually a type in the
static semantics and a value in the dynamic semantics. Such a sentence is
read ``against the background $A$, $\mathit{phrase}$ elaborates (or
evaluates) to $A'$''. The rules are formal; the semantic objects are the
subject of a small mathematical theory, given in the chapters that
introduce them, and the properties of the objects on which the rules rely
are stated there without proof.

Some notation is defined once and used throughout. Finite subsets, finite
maps and the modification $f+g$ of one finite map by another are defined in
Section~\ref{sec:static-core-compound}, together with the notation
$x^{(k)}$ for a sequence of length $k$; projection, injection and
modification of tuples of objects in
Section~\ref{sec:static-core-projection}; and the first and second options
$\opt{\ }$ and $\optt{\ }$ of the rules in
Section~\ref{sec:static-core-rules}. The conventions of the grammar figures
are stated with the grammar of the Core (Chapter~\ref{ch:syntax-core}).
Throughout, a metavariable decorated with primes or subscripts ranges over
the same class as the undecorated one.
